% chapters/design/exam_lifecycle.tex
% Sección 4.6: Ciclo de vida del examen (~1-1.5 páginas)
% ============================================================================
%
% GUÍA: Describir la máquina de estados de las instancias de examen.
%       Referencia a \ref{rf:lifecycle}.
%
%   --- DIAGRAMA 9: Diagrama de estados del examen ---
%   Mostrar todos los estados y transiciones. Estados probables:
%     - created       → Instancia creada tras OCR exitoso
%     - assigned      → Correctores asignados
%     - in_correction → Al menos un corrector ha empezado
%     - corrected     → Corrección completada
%     - published     → Notas publicadas, se abre ventana de revisión
%     - in_review     → Hay solicitudes de revisión pendientes
%     - final         → Notas definitivas
%     - archived      → Curso archivado
%
%   Transiciones con condiciones (quién puede disparar cada una,
%   qué precondiciones se deben cumplir).
%
%   \begin{figure}{fig:exam_states}{Diagrama de estados del ciclo de vida de una instancia de examen}
%     \image{}{}{exam_state_diagram}
%   \end{figure}
%
%   TODO: Crear la imagen figures/exam_state_diagram.png (o .pdf).
%         Herramientas: draw.io, PlantUML (statechart), Mermaid.
%         NOTA: Este diagrama es MUY importante para el tribunal; hazlo
%         claro y con buena resolución.
%
%   --- CONTENIDO TEXTUAL ---
%   1. Describir cada estado y su significado funcional.
%   2. Describir cada transición:
%      - Evento que la dispara.
%      - Actor que la puede ejecutar (rol).
%      - Precondiciones.
%      - Acciones asociadas (notificaciones, auditoría, etc.).
%   3. Explicar la implementación técnica:
%      - Campo de estado en el modelo ExamInstance.
%      - Validación de transiciones (patrón State o condicional).
%      - Concurrencia optimista (version field) para evitar conflictos.
%   4. Referencia a las notificaciones asociadas (\ref{rf:notifications}).
% ============================================================================

% TODO: Redactar la sección del ciclo de vida y crear el diagrama de estados.
